\section*{Chapitre \ref{ch:b2:continuous-reactors} --- Réacteurs continus et procédés industriels}

\begin{solution}{exo:b2:continuous-reactors:1}
$\tau = V/Q_v = 2000/0.50 = \qty{4.0e3}{s}$, soit \qty{1.1}{h}.
\end{solution}

\begin{solution}{exo:b2:continuous-reactors:2}
$k\tau = 1.2 \times 10^{-3} \times 1800 = 2.16$. Cuve~: $X = 2.16/3.16 = 0.68$~;
tube~: $X = 1 - \mathrm e^{-2.16} = 0.88$.
\end{solution}

\begin{solution}{exo:b2:continuous-reactors:3}
Tube~: $\tau = \ln 100/0.020 = \qty{230}{s}$. Cuve~: $X/(1 - X) = 99 = k\tau$,
$\tau = \qty{4950}{s}$, plus de vingt fois plus long~: aux \omterm{def:b2:continuous-reactors:conversion}{conversions} élevées, la
cuve agitée paie cher le fait de travailler à la concentration de sortie.
\end{solution}

\begin{solution}{exo:b2:continuous-reactors:4}
Par litre~: $120 \times 0.50 = \qty{60}{kJ}$ libérés, qui chauffent
\qty{4.2}{kJ/K}~: $\Delta T_{\text{ad}} = \qty{14}{K}$. Une panne du refroidissement
réchaufferait sensiblement le mélange, mais sans danger, sauf si le solvant est
proche de son point d'ébullition ou si une seconde réaction démarre dans ce domaine.
\end{solution}

\begin{solution}{exo:b2:continuous-reactors:5}
Cuve~: $Da = X/(1 - X)^2 = 0.8/0.04 = 20$, $\tau = 20/0.010 = \qty{2000}{s}$.
Tube~: $Da = X/(1 - X) = 4$, $\tau = \qty{400}{s}$. Rapport 5, contre $4/\ln 5 =
2.5$ pour l'ordre 1~: plus l'ordre est élevé, plus la cuve agitée pâtit
de sa faible concentration.
\end{solution}

\begin{solution}{exo:b2:continuous-reactors:6}
$k\tau = 5.0$. Trois cuves~: $X = 1 - (1 + 5/3)^{-3} = 0.947$. Une cuve~: $5/6 =
0.833$. Tube~: $1 - \mathrm e^{-5} = 0.993$.
\end{solution}

\begin{solution}{exo:b2:continuous-reactors:7}
$X = 0.70$ signifie $k\tau = 0.70/0.30 = 2.33$. Doubler le débit divise $\tau$ par deux~:
$k\tau = 1.17$, $X = 0.54$. Pour retrouver 70\,\%, il faut doubler le \omterm{def:b2:continuous-reactors:residence-time}{temps de passage}, donc le
volume.
\end{solution}

\begin{solution}{exo:b2:continuous-reactors:8}
Chaleur libérée~: $80 \times 2.0 \times 1.0 \times 0.90 = \qty{144}{kW}$. Emportée par le
flux de sortie, chauffé de \qty{330}{K} à \qty{350}{K}~: $4.0 \times 1.0 \times 20 =
\qty{80}{kW}$. La double enveloppe doit évacuer la différence, \qty{64}{kW}.
\end{solution}

\begin{solution}{exo:b2:continuous-reactors:9}
Dimensions linéaires $\times 10$~: volume $\times 1000$, aire de la paroi $\times 100$, aire
par unité de volume $\div 10$. La chaleur libérée croît comme le volume, la chaleur
évacuée par la paroi comme son aire~: le grand réacteur évacue, par litre,
dix fois moins de chaleur pour le même écart de température.
\end{solution}

\begin{solution}{exo:b2:continuous-reactors:10}
Tube~: $\tau = \ln(0.05/0.10)/(0.05 - 0.10) = \qty{13.9}{min}$, $c_B/c_{A,0} =
\frac{0.10}{-0.05}(\mathrm e^{-1.386} - \mathrm e^{-0.693}) = 0.50$. Cuve~: $\tau =
1/\sqrt{0.005} = \qty{14.1}{min}$, $c_B/c_{A,0} = 1.414/(2.414 \times 1.707) =
0.34$. Dans la cuve, du A frais et du B déjà formé sont mélangés dans le même volume~: une partie
de B séjourne toujours assez longtemps pour être convertie en C alors qu'une partie de A vient à peine
d'arriver~; dans le tube, chaque tranche a la même histoire.
\end{solution}

\begin{solution}{exo:b2:continuous-reactors:11}
Élever la température du fluide de refroidissement (et de l'alimentation) déplace la droite d'évacuation vers la
droite. L'intersection froide monte lentement jusqu'à ce que la droite devienne tangente au
coude inférieur du S~; au-delà, l'état froid disparaît et la cuve saute
sur la branche chaude (allumage). Quand on abaisse de nouveau la température, la cuve reste
sur la branche chaude jusqu'à ce que la droite devienne tangente au coude supérieur, et ne
retombe qu'alors (extinction), à une température plus basse que l'allumage~: un
cycle d'hystérésis.
\end{solution}

\begin{solution}{exo:b2:continuous-reactors:12}
\omterm{def:b2:continuous-reactors:reactors}{Réacteur discontinu}~: il reste 75\,\% des \qty{2.0}{mol/L}, $1.5 \times 100 = \qty{150}{kJ/L}$,
$\Delta T = 150/4.0 = \qty{37.5}{K}$ au-dessus de la température de travail, assez pour
atteindre un point d'ébullition ou une décomposition. Semi-continu~: au plus
\qty{0.10}{mol/L} n'a pas réagi, $\Delta T = 10/4.0 = \qty{2.5}{K}$. L'addition lente
plafonne l'énergie stockée dans le réacteur.
\end{solution}

\begin{solution}{pb:b2:continuous-reactors:1}
\textbf{1.} $r = k[\text{ester}][\ce{OH-}]$~; des alimentations égales et une
stœchiométrie 1 : 1 maintiennent les deux concentrations égales, $r = kc^2$.
\textbf{2.} $1/c - 1/c_0 = kt$ avec $c = 0.10c_0$~: $kc_0t = 9$, $t = 9/(0.11
\times 0.10) = \qty{818}{s}$, environ \qty{14}{min}.
\textbf{3.} $t_{1/2} = 1/(kc_0) = \qty{91}{s}$~; à l'ordre 2, la vitesse diminue comme
le carré de la concentration~: le temps nécessaire pour la diviser par deux dépend donc du point
de départ.
\textbf{4.} $Da = kc_0\tau = 0.011\,\tau$ ($\tau$ en secondes).
\textbf{5.} $Q_vc_0 - Q_vc - kc^2V = 0$, c'est-à-dire $c_0 - c = k\tau c^2$.
\textbf{6.} Avec $c = c_0(1 - X)$~: $c_0X = k\tau c_0^2(1 - X)^2$.
\textbf{7.} $Da = 0.90/0.10^2 = 90$.
\textbf{8.} $\tau = 90/0.011 = \qty{8.2e3}{s}$, \qty{2.3}{h}.
\textbf{9.} $V = 0.50 \times 8182 = \qty{4.1e3}{L}$, \qty{4.1}{m^3}.
\textbf{10.} La concentration de sortie, $0.010\,\unit{mol/L}$~: la vitesse est
partout cent fois plus faible qu'à l'entrée.
\textbf{11.} $Da = X/(1 - X) = 9$, $\tau = \qty{818}{s}$, $V = \qty{409}{L}$.
\textbf{12.} Dix fois moins de volume pour le tube.
\textbf{13.} $c_{j-1} - c_j = k\tau_1c_j^2$, où $\tau_1$ est le \omterm{def:b2:continuous-reactors:residence-time}{temps de passage}
d'une cuve.
\textbf{14.} $\tau_1 = 850/(3 \times 0.50) = \qty{567}{s}$, $k\tau_1 =
\qty{62.4}{L/mol}$. En résolvant $62.4c_j^2 + c_j - c_{j-1} = 0$~: $c_1 = 0.0328$,
$c_2 = 0.0163$, $c_3 = \qty{0.0100}{mol/L}$~: 90\,\% de \omterm{def:b2:continuous-reactors:conversion}{conversion}.
\textbf{15.} Le tube, ou la cascade (dix et cinq fois plus petits qu'une seule
cuve). Une cuve unique peut tout de même l'emporter par sa température uniforme, la facilité
de son nettoyage et de sa conduite, et parce qu'ici le volume coûte peu.
\textbf{16.} $\Delta T_{\text{ad}} = 55\,000 \times 100/(997 \times 4180) =
\qty{1.3}{K}$ (\qty{1.2}{K} à 90\,\%).
\textbf{17.} $55 \times 0.50 \times 0.10 \times 0.90 = \qty{2.5}{kW}$.
\textbf{18.} Non~: le flux de sortie emporte la chaleur avec une élévation d'environ
un kelvin~; aucun emballement n'est possible.
\textbf{19.} Vingt fois plus concentré~: l'élévation adiabatique devient
\qty{26}{K} et la chaleur \qty{50}{kW}, ce qui justifie un serpentin de refroidissement~; la vitesse $kc_0$
étant vingt fois plus élevée, la cuve pour 90\,\% serait vingt fois plus petite
(\qty{0.20}{m^3}).
\textbf{20.} $0.50 \times 0.10 \times 0.90 = \qty{0.045}{mol/s}$, $\times 86\,400
\times 82.03 = \qty{319}{kg}$ d'éthanoate de sodium par jour.
\textbf{21.} La constante de vitesse croît avec la température (loi d'Arrhenius)~: le
volume diminue~; le prix à payer est le chauffage de l'alimentation et, pour d'autres
réactions, une \omterm{def:b2:continuous-reactors:selectivity}{sélectivité} plus faible et davantage de vapeur.
\textbf{22.} La cuve agitée unique doit contenir $\boldsymbol{V \approx
\qty{4.1}{m^3}}$.
\end{solution}
